% preamble/geometry.tex
% 数学物理方法笔记 —— 页面几何与字体

% ── 页面几何 ──
\geometry{
	textheight=9in,
	textwidth=5.5in,
	top=1in,
	headheight=12pt,
	headsep=25pt,
	footskip=30pt
}

% ── 全局正文字体（英文 + 中文）──
\setmainfont{Times New Roman}    % 英文正文
\setCJKmainfont{SimSun}[AutoFakeBold={3}]          % 中文正文：宋体（伪粗体）

% ── 定理环境中文字体：仿宋 ──
\setCJKfamilyfont{theoremCJK}{FangSong}[AutoFakeBold={2.5}]   % 仿宋（伪粗体）
\newcommand{\theoremFont}{\CJKfamily{theoremCJK}}   % 只切换中文为仿宋

% ── 证明环境中文字体：楷体 ──
\setCJKfamilyfont{proofCJK}{KaiTi}[AutoFakeBold={2.5}]        % 楷体（伪粗体）
\newcommand{\proofFont}{\CJKfamily{proofCJK}}       % 只切换中文为楷体

% 主色调
\definecolor{mainblue}{HTML}{0E7882}
\definecolor{mediumblue}{HTML}{129BA8}
% 辅助色
\definecolor{lightblue}{HTML}{67D8E2}
\definecolor{grayblue}{HTML}{5C8A8E}
% 强调色
\definecolor{brightblue}{HTML}{F2C94C}
\definecolor{warningyellow}{HTML}{FDE047}
% 背景色
\definecolor{paleblue}{HTML}{F0FBFB}

\geometry{textheight=9in,textwidth=5.5in,top=1in}
